\section{Conclusion}


We have presented a hybrid post-quantum cryptographic architecture for blockchain
systems that combines BLS12-381 with ML-DSA-65 for signatures and X25519 with
ML-KEM-768 for key exchange. The architecture is provably at least as secure as the
stronger of the two components, adds acceptable overhead (1.43\,ms signing, 248
bytes/epoch), and supports a graceful migration path.

The key insight is that hybrid deployment is not a hedge---it is \emph{strictly
dominant}. During the transition period, we do not know whether lattice assumptions
or elliptic curve assumptions are stronger. The hybrid construction is secure
regardless of which assumption ultimately fails. The cost is modest. The benefit
is existential. Immediate hybrid deployment is the only rational choice.

%% ========================================================================
%%  BIBLIOGRAPHY
%% ========================================================================
\begin{thebibliography}{20}

\bibitem{fips203}
NIST, ``Module-Lattice-Based Key-Encapsulation Mechanism Standard (FIPS 203),''
National Institute of Standards and Technology, August 2024.

\bibitem{fips204}
NIST, ``Module-Lattice-Based Digital Signature Standard (FIPS 204),''
National Institute of Standards and Technology, August 2024.

\bibitem{fips205}
NIST, ``Stateless Hash-Based Digital Signature Standard (FIPS 205),''
National Institute of Standards and Technology, August 2024.

\bibitem{boneh2001bls}
D. Boneh, B. Lynn, and H. Shacham, ``Short Signatures from the Weil Pairing,''
\emph{Journal of Cryptology}, vol.~17, no.~4, pp.~297--319, 2004.

\bibitem{shor1997}
P.~W. Shor, ``Polynomial-Time Algorithms for Prime Factorization and Discrete
Logarithms on a Quantum Computer,'' \emph{SIAM Journal on Computing}, vol.~26,
no.~5, pp.~1484--1509, 1997.

\bibitem{stebila2020hybrid}
D. Stebila, S. Fluhrer, and S. Gueron, ``Hybrid Key Exchange in TLS~1.3,''
Internet-Draft, draft-ietf-tls-hybrid-design, 2020.

\bibitem{rfc9180}
R. Barnes, K. Bhargavan, B. Lipp, and C. Wood, ``Hybrid Public Key Encryption,''
RFC 9180, IETF, February 2022.

\bibitem{bindel2019hybrid}
N. Bindel, J. Brendel, M. Fischlin, B. Goncalves, and D. Stebila, ``Hybrid Key
Encapsulation Mechanisms and Authenticated Key Exchange,'' \emph{Post-Quantum
Cryptography (PQCrypto)}, pp.~206--226, Springer, 2019.

\bibitem{qrl2018}
P. Waterland et al., ``The Quantum Resistant Ledger,'' QRL Foundation, 2018.

\bibitem{eip7562}
Ethereum Foundation, ``EIP-7562: Account Abstraction Validation Scope Rules,''
Ethereum Improvement Proposals, 2023.

\bibitem{grover1996}
L.~K. Grover, ``A Fast Quantum Mechanical Algorithm for Database Search,''
\emph{Proceedings of the 28th Annual ACM Symposium on Theory of Computing},
pp.~212--219, 1996.

\bibitem{mosca2018}
M.~Mosca, ``Cybersecurity in an Era with Quantum Computers: Will We Be Ready?''
\emph{IEEE Security \& Privacy}, vol.~16, no.~5, pp.~38--41, 2018.

\end{thebibliography}

